\section{Action, conventions, and scalar derivatives}\label{app:action}
We use metric signature $(-,+,+,+)$ and the definitions of Sec.~\ref{sec:theory}. The scalar--vector action is Eq.~\eqref{eq:aest-action}, with the nonseparable function Eq.~\eqref{eq:fullF}. On the homogeneous branch $\Y=0$,
\begin{align}
F_{,\Y}&=\Gamma(\Q),\\
F_{,\Y\Q}&=\Gamma_{,\Q},\\
F_{,\Q}&=\mathcal F_{T,\Q},\\
F_{,\Q\Q}&=\mathcal F_{T,\Q\Q}.
\end{align}
At the static screened point $\Q=\Q_0$,
\begin{equation}
\Gamma(\Q_0)=0,
\end{equation}
while $F_{,\Y\Q}$ need not vanish. Since $\Y$ is quadratic in spatial scalar gradients, the first nonzero mixed $F_{,\Y\Q}$ contribution is cubic in perturbations, explaining why the new coefficient enters the cubic cutoff calculation but not as an independent direct quadratic source.

For completeness, varying the Type-II auxiliary multiplier in Eq.~\eqref{eq:vcdm-aux} gives
\begin{equation}
-\frac32\lambda_C-(K+\varphi_V)=0,
\end{equation}
which yields the reduced term Eq.~\eqref{eq:vcdm-red}. In the scalar perturbation sector at nonzero wavenumber, the auxiliary equations set $\delta\varphi_V=0$ and fix the associated multipliers algebraically. Thus the acceleration sector alters the constraint equations and background reconstruction without adding a local scalar oscillator. This is the same degree-of-freedom property that characterizes VCDM in the Type-II minimally modified gravity literature \cite{DeFeliceMaedaMukohyamaPookkillath2022,GanzMartensMukohyamaNamba2023}.

The full scalar Euler--Lagrange equation follows from the shift symmetry. Writing the scalar contribution as $\sqrt{-g}\,L(\nabla\phi,A,g)$, it takes current-conservation form $\nabla_\mu \mathcal J^\mu_\phi=0$, where the timelike projection contains $F_{,\Q}$ and the projected spatial part contains $F_{,\Y}q^{\mu\nu}\nabla_\nu\phi$. In FLRW the spatial contribution vanishes identically, leaving the conserved quantity derived in Appendix~\ref{app:flrw}. This derivation is also the reason $\Gamma(\Q)$ cannot be adjusted independently once the current solution and temporal function are fixed.
